% Transcription — fonds Alexandre Grothendieck, shelfmark 43, batch 6 (pages 101-120).
% Established from the facsimile archives/batches/43/batch-06.pdf (PDF page k =
% archivists' page 100+k, no cover sheet), cut from a copy of 43.pdf fetched
% through the project's own relay
% (https://grothendieck-relay.onrender.com/source/43.pdf); tiles in
% archives/tiles/43/p101-p120. Per the skill transcribe-grothendieck. The folder
% is 207 pages in eleven batches, transcribed in parallel. Batches 1, 2, 5 and
% 9-11 were read for the folder's conventions; batch 5 (committed during this
% pass) ends on the work plan 1)-8) of p. 99, and p. 101 here continues that
% same plan with items 9)-14), so the notation of batch 5 is followed:
% \underline{\mathrm{Ext}} for his underlined Ext (doubly underlined where he
% doubles the stroke), \mathrm{Ext}^{i}_{k\text{-}\mathrm{gr}}, \mathbb{G}_m,
% \mathcal{C} for his curly C, K the fraction field of the complete DVR V with
% residue field k, J_{K/k}, T and S the functors of the diagram of p. 105.
%
% Pass: Opus 5.5 (claude-opus-5-5), 2026-09-26 — first pass, unchecked against the pages by a human.
%
% Dating: the folder has its own inventory date, carried verbatim in \dating{}
% with no group sentence (the skill adds the group « Jacobiennes » (43-44)
% only for undated folders). No page of the notes carries a date; the
% cancelled English to-do list upside down at the foot of p. 117 names
% « September 1963 » (said in a \note there, the list itself not transcribed).
%
% Copies. No page of this batch can be identified as a photocopy. All the
% notes are pencil or blue / blue-black ink on original paper; several leaves
% show the typing of their other side through the paper (108 shows 109 in
% mirror, 113 shows 114, 119 shows 120), which a photocopy would not.
%
% Pages skipped: 102 and 104, two typed leaves of an English paper on Stiefel
% manifolds and « intrinsic join » (typed p. 21 and the REFERENCES), scanned
% upside down, with nothing in his hand — the same paper as batch 5's skipped
% p. 100; 110, a blank sheet.
%
% Typescripts. Pages 114 and 120 are leaves of a typescript of EGA III (typed
% pages « III-49 » and « III-42 »), with ink corrections of a hand not
% established here and yellow highlighting: one French \note each, the ink
% interventions listed, the typed text not re-set. Page 109 is a typed leaf
% headed « IV 7 Oubli au nº 7.6. » (Proposition 6.7, Lemme 6.8, EGA IV
% numbering) with heavy ink revision and an ink N.B. in his hand, cancelled by
% three long strokes; whether this text was published in that form was not
% established, so, being short, it is given entire (typed text plain, his ink
% as \add / \struck / \marginal), following folder 27 batch 1.
%
% Pages summarised: none. Pages 107 (upside down, read turned) and 112
% (cancelled by two diagonal strokes) are fast prose and carry a heavy
% apparatus; the scattered formula sheets 103, 106, 108, 111, 113, 115, 116,
% 118, 119 are given block by block, their layout said in \note{}s; the
% diagram of p. 105 is redrawn as a tikz-cd grid.
%
% His pagination: none on his own notes. Typed page numbers III-42, III-49 on
% the EGA III leaves. Numbered runs: p. 101 continues the plan 1)-8) of p. 99
% (batch 5) with 9) Application 1 (the isomorphism H^i(K, gamma_K) ->
% Ext^i(J_{K/k}, gamma), « idem aux corps de classes de Serre »), 10)
% Application 2 (autoduality, Witt vectors), 11) Application 3 (cohomology of
% abelian varieties), 12) L'isom. de Lang, 13) Application 4, Dualité de Tate,
% 14) Compléments aux égales caractéristiques / Théorèmes de Rosenlicht
% locaux. Questions (i)-(iii) on T° and T* (p. 103). Cases 1°) G fini
% ordinaire (p. 106), 2°) G = mu_p, 3°) G = alpha_p, 3°) G = G_m (p. 108 —
% the 3° is used twice). (1)-(3) in the box of p. 116. The EGA IV plan
% 1.-15. of p. 117. (i)-(iii) and 1°)-2°) on p. 119. Corollaire ?.9.10 (p. 107,
% the first digit unreadable), Proposition 6.7 and Lemme 6.8 (p. 109),
% Remarque 2.7 and Prop. 2.8 (p. 115, upside down). No internal page reference.
%
% What the batch is about. Duality for commutative algebraic groups over a
% local field K = Frac V with residue field k, continuing batch 5. P. 101: the
% rest of the plan, and the question H^i(K, gamma_K) ≃ Ext^i_{k-gr}(J_{K/k},
% gamma) for i > 0, with a bracketed genealogy (Lang, Serre, Tate). P. 103: the
% categories of finite complexes of K-groups and k-quasi-groups, T : KGr_K ->
% KQuGr_k, Gamma*_K = Gamma*_k T, questions on T°, and the spectral sequences
% H^p(k, T^q(G)) => H*(K, G), H^p(k, Ext^q_{K/k}) => Ext_{K-gr}. P. 105: the
% square of dualities D (Cartier), Delta (Serre), Theta (Pontrjagin) linked by
% T (« théorie de dualité géom. ») and S (« Frobenius-Lang ») for k finite, and
% Ext^i_{k-gr}(G, Z/nZ) ≃ Ext^i(SG, Z/nZ(-1)) with the check on Z/nZ. P. 106:
% the pairing H^i(A) x H^{2-i}(B) -> H^2(K, G_m), Ext^i_{k-gr}(G, Z/nZ) ≃
% Hom(H^{1-i}(K,G), Z/nZ) and the same with Q/Z, case of finite G. P. 107:
% Cor. ?.9.10, semi-continuity of a number of components pi(z) for f of finite
% type, reduced to a DVR. P. 108: Ext^i_{k-gr}(G, Q/Z) for G = mu_p, alpha_p,
% G_m, with H^i(k, G_m). P. 109: the EGA IV typescript (semi-continuity of
% total geometric multiplicities). P. 111: a table of mu_n, Z/nZ, mu_p, Z/pZ,
% alpha_p, G_m, Z with their gerbes (Grbg) and H^0, H^1 over K. P. 112
% (cancelled): reciprocity homomorphisms phi_x : Gal(K-bar_x/K_x) -> Gal(K'/K),
% decomposition and inertia groups, K_X^D, K_X^I. P. 113: H^i(K, D(G)) ≃
% Ext^{1-i}(f_*G, Q/Z), f_* D(G) ≃ (T f_*G)(1), and the spectral sequence
% Ext^p(H^{-q}(K, G), Q/Z) => H*(K, DG). P. 115: H*(Delta^2 T(A)) <=
% (H^p Delta^2)(H^q(K, A)), Ext^0(P, Q/Z) = 0. P. 116: for abelian varieties
% A, B dual: H^1(K,A) ≃ Ext^1_{k-gr}(B, Q/Z), pi_0(A), pi_0(B) dual, 0 -> U
% -> A -> A' -> 0 and mu(A) = 2 alpha(A°) + mu(A°). P. 117: plan of EGA IV,
% 1.-15. Pp. 118-119: Ext^i(A, mu_infinity), Ext^i(A, _n mu) ≃ H^{i-1}(K, _nB)
% ≃ Ext^{i-1}(T(A), Z/nZ), H^i(K, A), the exact sequences with DH^1(K, A),
% and the divisibility of H^1(K, A)'.
%
% Readings to check first: the names in item 10) of p. 101 and the struck
% Hom line there; the right-hand bracketed glosses of p. 103; the legend of
% p. 105 (« Greenberg »?); the whole of pp. 107 and 112; the N.B. of p. 109;
% the glosses under mu(A) on p. 116; item 8. and 14. of p. 117; the small
% superscript marks after A(K) and H^1(K, _infinity B) on p. 118.

\documentclass[11pt,a4paper]{article}
% Le préambule commun à toutes les transcriptions du fonds Grothendieck.
%
% Il tient en une idée : **l'incertitude se compose, elle ne se résout pas.**
% Une transcription qui lisse un mot douteux a détruit la seule chose pour
% laquelle on la faisait — un lecteur doit pouvoir savoir, à chaque endroit,
% ce qui était sur le feuillet et ce qui a été ajouté. Les six macros qui
% suivent sont donc l'appareil critique, et non de la décoration.
%
% Les mêmes commandes sont reconnues par `scripts/render.mjs`, qui produit la
% vue de lecture du site. Une macro ajoutée ici doit l'être aussi là-bas,
% faute de quoi le rendu échouera bruyamment — ce qui est voulu : un rendu
% silencieusement faux est pire qu'un rendu absent.

\ProvidesPackage{grothendieck}[2026/08/08 Transcriptions du fonds Grothendieck]

\usepackage{amsmath,amssymb,amsthm}
\usepackage{mathrsfs}
\usepackage{stmaryrd}
% Les diagrammes commutatifs sont partout dans ce fonds : c'est la langue
% même de la géométrie algébrique de Grothendieck, et une transcription qui
% les rendrait en prose ne transcrirait rien.
\usepackage{tikz-cd}
% Français par défaut — c'est la langue des feuillets — mais l'anglais est
% chargé aussi : la traduction et le résumé sont des documents anglais, et la
% typographie française (espaces avant les deux-points) y serait une faute.
% Ces documents basculent par \selectlanguage{english} en tête de corps.
\usepackage[english,french]{babel}
\usepackage{fontspec}
\usepackage{geometry}
\geometry{a4paper,margin=25mm,marginparwidth=28mm}
\usepackage{marginnote}
\usepackage{xcolor}
% ulem, not soul. `soul`'s \st and \ul are built on a character-by-character
% scan that XeLaTeX's font machinery breaks: the document fails with an
% undefined control sequence rather than degrading. ulem does the same two jobs
% and is Unicode-engine safe. `normalem` stops it hijacking \emph, which the
% transcriptions use for Grothendieck's own emphasis.
\usepackage[normalem]{ulem}
\usepackage{hyperref}
\hypersetup{colorlinks=true,linkcolor=black,urlcolor=black,pdfborder={0 0 0}}

% --- Métadonnées du lot -----------------------------------------------------
% Reprises telles quelles par le script de rendu, qui les affiche en tête de
% la vue de lecture. Elles fixent ce que le fichier prétend couvrir : sans
% elles, une transcription flotte sans point d'attache dans un fonds de seize
% mille feuillets.

\newcommand{\folder}[1]{\gdef\grothFolder{#1}}
\newcommand{\batch}[1]{\gdef\grothBatch{#1}}
\newcommand{\leaves}[2]{\gdef\grothFirst{#1}\gdef\grothLast{#2}}
\newcommand{\foldertitle}[1]{\gdef\grothFoldertitle{#1}}
\gdef\grothFolder{?}\gdef\grothBatch{?}\gdef\grothFirst{?}\gdef\grothLast{?}\gdef\grothFoldertitle{}

% --- Le résumé --------------------------------------------------------------
%
% Il ouvre la lecture modernisée, sous le titre « Résumé », et il est composé
% en petit. Ce n'est pas un ornement : c'est le seul passage du document qui
% ne soit pas des mathématiques, et un lecteur doit pouvoir le reconnaître
% avant de l'avoir lu — puis le sauter s'il connaît déjà le sujet, ou s'y
% arrêter s'il ne le connaît pas. Le filet qui le suit marque la reprise du
% corps.

\newenvironment{resume}
  {\small\setlength{\parskip}{0.45em}}
  {\par\medskip\hrule\medskip}

% --- Mots-clés ---------------------------------------------------------------
%
% En anglais, en clôture du résumé de la lecture modernisée : le vocabulaire
% moderne sous lequel chercher ce que les feuillets construisent. Le manifeste
% du site (`npm run manifest`) les extrait pour étiqueter la cote entière —
% cette ligne est la source unique des tags, il n'y a pas de fichier à côté.
\newcommand{\keywords}[1]{\par\smallskip\noindent
  {\footnotesize\textsc{Keywords} --- \textit{#1}}\par}

% --- L'appareil critique ----------------------------------------------------

% Le numéro de feuillet, en marge. C'est l'ancre qui permet de régler un
% désaccord sur une formule en regardant *un* feuillet plutôt qu'une chemise
% de six cents. Le site s'en sert aussi pour tourner les pages du fac-similé
% quand on fait défiler la transcription.
\newcommand{\leaf}[1]{%
  \marginnote{\footnotesize\color{gray!70}\textsf{#1}}%
  \label{leaf:#1}\ignorespaces}

% Illisible : on ne devine pas. Un mot inventé qui se lit comme les autres est
% le pire résultat possible.
\newcommand{\ill}{\textcolor{red!60!black}{[\,\ldots\,]}}

% Lecture douteuse : le mot est donné, et signalé comme incertain.
\newcommand{\uncertain}[1]{\uline{#1}}

% Ajout de l'éditeur — une lettre manquante, un mot sous-entendu.
\newcommand{\add}[1]{\textcolor{blue!50!black}{[#1]}}

% Biffé par Grothendieck lui-même. Ce qu'il a rayé fait partie du document :
% une rature dit souvent où la pensée a changé de direction.
\newcommand{\struck}[1]{\sout{#1}}

% Note du transcripteur, sur l'état matériel du feuillet.
\newcommand{\note}[1]{\par\smallskip{\small\color{gray!60!black}\textit{[#1]}}\par\smallskip}

% Note marginale de Grothendieck, distincte du corps du texte.
\newcommand{\marginal}[1]{\par\smallskip\noindent\hspace{1em}%
  {\small\color{gray!60!black}#1}\par\smallskip}

% --- Titre ------------------------------------------------------------------

\AtBeginDocument{%
  \begin{center}
    {\large\bfseries Fonds Alexandre Grothendieck --- cote n\textsuperscript{o}~\grothFolder}\\[2pt]
    {\itshape\grothFoldertitle}\\[4pt]
    {\small feuillets \grothFirst--\grothLast\ (lot \grothBatch)}\\[2pt]
    {\footnotesize\color{gray!60!black}Transcription établie d'après le fac-similé numérisé
     par l'Université de Montpellier.\\
     Les lectures incertaines sont soulignées, les illisibles notées [\,\ldots\,],
     les ajouts entre crochets.}
  \end{center}
  \vspace{1em}}

\endinput


\folder{43}
\batch{6}
\pages{101}{120}
\dating{[à partir de 1961-vers 1968]}
\watermark{Édition de démonstration}
\foldertitle{Dualité des groupes. Jacobiennes généralisées etc : notes et copies de notes manuscrites (s.d.)}

\begin{document}

\page{101}
\note{au crayon. La numérotation 9)--14) continue le plan de travail
1)--8) de la page 99 (batch 5).}
\begin{enumerate}
\item[9)] \emph{Application 1} Le isom.\
$H^{i}(K, \gamma_{K}) \xrightarrow{\ \sim\ } \mathrm{Ext}^{i}_{K\text{-}\mathrm{gr}}(J_{K/k}, \gamma)$

\struck{et les corps de classes de Serre}

(Idem aux corps de classes de Serre)
\note{à la fin de la ligne, l'indice de $J$ et le second argument sont
coupés par le bord du feuillet ; l'indice du $\mathrm{Ext}$ se lit
« $K$-gr » et non « $k$-gr », lecture incertaine.}
\item[10)] \emph{Application 2} Autodualité des \struck{\ill{}}
\uncertain{$H W_{n}(K)/$}\ill{} \ill{}
\marginal{?}
\note{le chiffre 10 est repassé ; en marge gauche, un « ? » et un double
trait oblique. La fin de la ligne, serrée contre le bord, porte
$W_{n}(K)$ et, en dessous, peut-être « Witt » ; lecture très incertaine.}
\item[11)] \emph{Application 3} Cohomologie des V.A.
\item[12)] L'isom de Lang \struck{pour un \ill{}} (dans un corps $k$ fini)
\item[13)] \emph{Application 4} Dualité de Tate.
\item[14)] Compléments aux égales caractéristiques
\struck{\ill{}}

\emph{Théorèmes de Rosenlicht locaux}
\note{la première ligne du point 14) est barrée d'un trait fin ; « Théorèmes
de Rosenlicht locaux », souligné, est écrit dessous, sans doute en
remplacement.}
\end{enumerate}
\marginal{?}
\[
H^{i}(K, \gamma_{K}) \simeq \mathrm{Ext}^{i}_{k\text{-}\mathrm{gr}}(J_{K/k}, \gamma)
\]
pour $i > 0$, tt $k$-groupe $\gamma$, \uncertain{grâce} à \ill{}
canonique de $J(K)/J(V)'$.

\struck{$\mu_{p} \otimes \mu_{p}$ \quad
$\mathrm{Hom}(\mu_{p} \otimes \mu_{p}, \mathbb{G}_{m}) = \mathrm{Hom}(\mu_{p}, \mathbb{Z}/p\mathbb{Z})$}
\note{deux lignes au crayon en bas à gauche, barrées d'un trait ondulé ; la
lecture du second membre est incertaine.}

\begin{quote}
\emph{Lang} : géométrisation du corps de classes global par Rosenlicht
+ Lang

\emph{Serre} : géométrisation des corps de classes local

\emph{Tate} : Dualité de la cohomologie des V.A.\ sur les corps
\uncertain{$p$-adiques}.
\end{quote}
\note{les trois lignes sont réunies à gauche par un grand crochet repassé.}

\page{103}
\note{au crayon. En haut à gauche, sous $k \leftarrow V \to K$, un mot
entouré suivi de « ? », barré de plusieurs traits (\struck{\ill{}}).}
\[
k \longleftarrow V \longrightarrow K
\]

\begin{tikzcd}
K\,\mathrm{Gr}_{K} \arrow[d, "T"'] \\
K\,\mathrm{QuGr}_{k}
\end{tikzcd}

\note{à côté de la flèche $T$, un $S$ très noirci, sans doute biffé. En face
de $K\,\mathrm{Gr}_{K}$ : « complexes finis de \struck{\ill{}} $K$-groupes
\struck{\ill{}} » et, à droite, « [Coh.\ de $K$ ne se calcule pas
galoisiennement] » ; en face de $K\,\mathrm{QuGr}_{k}$ : « complexes finis
de $k$-quasi-groupes » et « [Coh.\ de $k$ se calcule galoisiennement,
\ill{} on doit \uncertain{pouvoir} remplacer $k$ par une clôture
\uncertain{parfaite}\ldots] ».}
\[
T(G) = T^{\circ}\mathcal{C}(G)
\]
\[
\Gamma^{*}_{K} G \simeq \Gamma^{\circ}_{K}\mathcal{C}(G) \simeq \Gamma^{\circ}_{k} T^{\circ}\mathcal{C}(G) \simeq \Gamma^{*}_{k} T(G)
\]
\note{la lettre cursive $\mathcal{C}$ est celle de la page 105 (« groupes
top.\ abéliens loc.\ compacts ») ; dans le premier terme, l'indice de
$\Gamma$ est surchargé.}

\struck{Question}
\begin{enumerate}
\item[(i)] $T^{\circ}$ transforme injectifs en $\Gamma^{\circ}_{k}$-acycliques ??
\item[(ii)] $T^{\circ}$ de dim coh.\ $\leqslant 2$ [et même $\leqslant 1$
si on \struck{excepte} \add{exclut} les composantes du type $\mathbb{Z}$]
\item[(iii)] $T^{*}$ commute à \emph{« l'extension du corps résiduel $k$ »}.

i.e.\ est « de nature géométrique »,
\end{enumerate}
\[
\underline{\underline{H}}^{*}(K, G^{\cdot}) \Longleftarrow H^{p}\bigl(k, \underline{\underline{T}}^{q}(G^{\cdot})\bigr)
\qquad
\underline{\underline{H}}^{*}(K, G^{\cdot}) \simeq \underline{\underline{H}}^{*}\bigl(k, T(G)\bigr)
\]
\note{le $\simeq$ est écrit verticalement sous $\underline{\underline{H}}^{*}(K, G^{\cdot})$.}
\[
\Bigl[\ \underline{\underline{\mathrm{Ext}}}^{\cdot}_{K\text{-}\mathrm{gr}}(F, G) \Longleftarrow H^{p}\bigl(k, \underline{\mathrm{Ext}}^{q}_{K/k}(F, G)\bigr) \Longleftarrow H^{p}\bigl(K, \underline{\mathrm{Ext}}^{q}_{K}(F, G)\bigr)
\]
\note{les deux $\mathrm{Ext}^{q}$ de droite sont soulignés d'un trait droit
et d'un trait ondulé ; on ne rend que le premier.}

\page{105}
\note{à l'encre bleu-noir, puis au crayon. Le diagramme est redessiné ;
entre les colonnes, deux annotations obliques : « Théorie de dualité géom. »
dans le carré du haut, « Frobenius-Lang » dans celui du bas, chacune
accompagnée d'un « $\approx$ » tourné vers la colonne de droite.}

\begin{tikzcd}[column sep=large]
K^{-}_{\mathrm{fini}}(K)^{\circ} \arrow[r, "D"] \arrow[d, "T"'] & K^{+}_{\mathrm{fini}}(K) \arrow[d, "T(-1)"] \\
K^{-}(k)'^{\circ} \arrow[r, "\Delta"] \arrow[d, "S"'] & K^{+}(k)' \arrow[d, "S(-1)"] \\
K^{-}\mathcal{C}^{\circ} \arrow[r, "\Theta"] & K^{+}\mathcal{C}
\end{tikzcd}

\begin{quote}
$D$ dualité de Cartier

$\Delta$ dualité de Serre

$\Theta$ dualité de Pontrjagin

$\mathcal{C}$ groupes top.\ abéliens loc.\ compacts

(Le $'$ indique tot.\ discontinus qui n'ont pas de composante
$\mathbb{Z}$)

$k$ corps \emph{fini}

$T$ foncteurs \uncertain{Greenberg}-\ill{}-\ill{}

$S$ foncteurs $\Gamma_{k}$
\end{quote}
\note{« fini » est souligné trois fois. $D$, $\Delta$, $\Theta$ sont réunis
par une accolade, $T$ et $S$ par une autre ; le second nom de la légende de
$T$, sur deux lignes, n'est pas lu.}

Comment obtient-on le carré du bas ?

\struck{\ill{}}
\[
\mathrm{Ext}^{i}_{k\text{-}\mathrm{gr}}(G, \mathbb{Z}/n\mathbb{Z}) \simeq \mathrm{Ext}^{i-1}_{\mathrm{gr.top.}}(SG, \mathbb{Z}/n\mathbb{Z}) = \mathrm{Ext}^{i}\bigl(SG, \mathbb{Z}/n\mathbb{Z}(-1)\bigr)
\]
\note{le dernier membre est écrit sous l'avant-dernier, relié par un
« $=$ » vertical.}
\[
G \xrightarrow{\ \deg i\ } \mathbb{Z}/n\mathbb{Z}
\qquad
SG \xrightarrow{\ \deg i\ } S(\mathbb{Z}/n\mathbb{Z}) \xrightarrow{\ \deg 0\ } \mathbb{Z}/n\mathbb{Z}(-1)
\]
\[
\underline{\underline{\mathrm{Ext}}}^{1}_{\mathrm{gr.top}}\bigl(S(\mathbb{Z}/n\mathbb{Z})(-1), \mathbb{Z}/n\mathbb{Z}(-1)\bigr) = \underline{\underline{\mathrm{Ext}}}^{1}_{k\text{-}\mathrm{gr}}\bigl(S(\mathbb{Z}/n\mathbb{Z}), \mathbb{Z}/n\mathbb{Z}\bigr)
\]
\[
= \mathrm{Hom}\bigl(H^{0}S(\mathbb{Z}/n\mathbb{Z}), \mathbb{Z}/n\mathbb{Z}\bigr) \simeq \mathbb{Z}/n\mathbb{Z} \quad !
\]
\note{l'exposant du premier $\mathrm{Ext}$ est un petit signe peu lisible, lu $1$ sans certitude ;
le « $(-1)$ » du premier argument est surchargé. Au bas de la page, isolé
derrière un trait vertical : « Th. ».}

\page{106}
\note{au crayon.}
\[
\boxed{H^{i}(A) \times H^{2-i}(B) \longrightarrow H^{2}(K, \mathbb{G}_{m})}
\]
\[
H^{1}({}_{n}A) \longrightarrow
\qquad\qquad
\underline{\underline{\varinjlim}}\ H^{1}(K, A)
\]
\note{le « 1 » de $H^{1}({}_{n}A)$ est écrit sur un autre chiffre ; le
$\varinjlim$ souligné deux fois est placé au-dessus de $H^{1}(K, A)$.}
\[
\mathrm{Ext}^{i}_{k\text{-}\mathrm{gr}}(G, H) \simeq \mathrm{Ext}^{i-2}(\mathrm{R}\Gamma G, H) = \mathrm{Hom}(H^{2-i}\ldots
\]
\note{le premier $\mathrm{Ext}$ est précédé d'un signe biffé ; la ligne
s'arrête sur « $(H^{2-i}(\,$ ». Dessous, isolé : « $H^{i}($ ».}
\struck{$H^{i}$}
\[
\mathrm{Ext}^{i}_{k\text{-}\mathrm{gr}}(G, \mathbb{Z}/n\mathbb{Z}) \simeq \mathrm{Hom}\bigl(H^{1-i}(K, G), \mathbb{Z}/n\mathbb{Z}\bigr)
\]
\[
\mathrm{Ext}^{i}_{k\text{-}\mathrm{gr}}(G, \mathbb{Q}/\mathbb{Z}) \simeq \mathrm{Hom}\bigl(H^{1-i}(K, G), \mathbb{Q}/\mathbb{Z}\bigr)
\]
\note{dans la seconde ligne, $\mathbb{Q}/\mathbb{Z}$ est écrit par-dessus
$\mathbb{Z}/n\mathbb{Z}$, aux deux places ; l'exposant $1-i$ y est
surchargé.}

\struck{1) $G$ gpe fini}

\emph{1°)} \emph{$G$ gpe fini ordinaire} \quad $\mathbb{Z}/n\mathbb{Z}$

\struck{$H^{2-i}$ : $H^{i}(K, G) \times H^{2-i}(K, \hat{G})$}
\[
H^{i}(K, \mathbb{Z}/n\mathbb{Z}) \times H^{1-i}(K, \mathbb{Z}/n\mathbb{Z}) \longrightarrow \mathbb{Q}/\mathbb{Z}
\]
\[
\mathbb{Z}/n\mathbb{Z} \qquad \mathbb{Z}/n\mathbb{Z}
\]
\note{l'exposant $1-i$ de l'accouplement est une lecture incertaine (peut-être $2-i$). Les deux $\mathbb{Z}/n\mathbb{Z}$ de la dernière ligne sont écrits
sous les deux facteurs de l'accouplement.}

\page{107}
\note{page écrite tête-bêche par rapport au numéro des archivistes, lue
retournée ; au crayon, écriture rapide, très raturée. Tout le premier
paragraphe est bordé à gauche d'un trait vertical ; un long trait oblique
traverse la page de haut en bas ; une large tache d'encre couvre le milieu
de la page. Les formules passent, la prose souvent pas.}

\emph{Corollaire} \ill{}.9.10. Soit $f \colon X \to Y$ un morphisme
\struck{\ill{}} de type fini, \uncertain{propre} \uncertain{avec} $Y$ loc.\
noeth., $y \in Y$, \ill{} \uncertain{suppose} $f$ univ.\ ouvert en les pts de
$X_{y}$, et $X_{y}$ \uncertain{séparable}. \struck{\ill{} \ill{} \ill{}
fonction} Pour tout $z \in Z$, \uncertain{soit} $\pi(z) = $ \ill{} \ill{}
\ill{} \struck{\ill{}} \uncertain{composantes} \ill{} \ill{} de $X_{z}$. Alors
la fonction $z \mapsto \pi(z)$ est \struck{\ill{}} \add{continue}
\struck{\ill{}} au \uncertain{voisinage} de $y$. [\ill{} \ill{} continue en
$y$].
\note{la marque en haut à droite, au-dessus de « morphisme », est un mot
biffé ; « $\pi(z) = $ » est suivi d'une fin de ligne raturée où l'on devine
$\pi_{0}(\ldots)$ ; le mot après la seconde « fonction » est dans un cadre
raturé.}

\struck{\ill{} \ill{} dit \ill{} \ill{} \ill{} $\tfrac{1}{2}$ continue
supérieurement \ill{} \ill{} \ill{} « \ill{} continuité ». Il \ill{} \ill{}
\ill{} \ill{} \ill{} supérieurement \ill{} \ill{} (\ill{} \ill{} \ill{}
cont.)}
\note{passage encadré à gauche et barré de plusieurs traits obliques et
d'une longue rature horizontale.}

\ill{} \ill{} \ill{} \ill{} $y'$ de $y$, on a $\pi(y) \geqq \pi(y')$. Le
procédé habituel nous ramène alors au cas où $Y$ est le spectre d'un anneau
de valuation discrète, et où les composantes connexes de $X_{y}$ (l'indice est une lecture incertaine)
\ill{} géométriquement \ill{}. \ill{} \ill{} \ill{} supposer $X$ connexe,
$X_{1}$ \ill{} connexe, \struck{\ill{} fini} donc \struck{$\pi_{0}(X_{1})$}
$\pi(X_{1}) = 1$, \ill{} \ill{} comme $X_{0} \neq \emptyset$, \ill{} \ill{}
fini \ill{}
\note{à droite, deux croquis : un secteur coupé de deux arcs, et une figure
portant « $X_{1} \neq \emptyset$ », « $X = \emptyset$ » et une accolade
marquée $\alpha_{p}$ (?). Au bas, raturé et entrelacé de traits, un reste de
suite « $\ldots \to W \to W/N \to \ldots$ », et en marge droite, écrit en
long, « \uncertain{cf.} \ill{} V.A.\ \ill{} Lang ».}

\page{108}
\note{au crayon, sur un feuillet dont le verso est le dactylogramme de la
page 109, visible par transparence.}

\emph{2°)} $G = \mu_{p}$ \qquad
$\underline{\mathrm{Ext}}^{i}_{k\text{-}\mathrm{gr}}(G, \mathbb{Q}/\mathbb{Z}) = \lbrace 0,\ 0,\ 0,\ 0$
\note{les valeurs sont écrites en colonne, à droite d'une accolade ; le
second argument est écrit sur une autre lettre.}

\struck{$H^{0}(k, \mu)$} \quad $\mathbb{Z}/p\mathbb{Z}$ \quad $\mu_{p}$

\struck{$H^{1}$} \quad $K^{*}/K^{*p}$
\[
H^{1}(k, \mu_{p}) \qquad K^{*}/K^{*p}
\]
\[
\boxed{H^{i}(k, \mathbb{G}_{m}) = \qquad k^{*} \xrightarrow{\ p\ } k^{*}}
\]
\note{le cadre, arrondi, enferme ces deux termes ; le second membre de
l'égalité n'est pas écrit.}

\emph{3°)} $G = \alpha_{p}$ \quad $\times p$
\note{le « $\times p$ » est écrit sous $\alpha_{p}$.}

\emph{3°)} $G = \mathbb{G}_{m}$
\[
\underline{\mathrm{Ext}}^{i}(\mathbb{G}_{m}, \mathbb{Q}/\mathbb{Z}) = \lbrace 0,\ (\text{racines de l'unité})^{\vee},\ 0,\ 0
\]
\[
\underline{\underline{\mathrm{Ext}}}^{i}_{k\text{-}\mathrm{gr}}(\mathbb{G}_{m}, \mathbb{Q}/\mathbb{Z}) = H^{i-1}(k, \text{racines de l'unité}) \quad \lbrace \text{racines de l'unité},\ 0,\ 0,\ 0
\]
\note{les valeurs sont écrites en colonne à droite d'accolades ;
« racines de l'unité » de la seconde colonne est plié sur trois lignes. Le
« ; » entre les deux arguments est de lui.}
\[
H^{i}(k, \mathbb{G}_{m}) = \lbrace k^{*},\ 0,\ 0,\ 0
\]
\struck{$k^{*}$}
\[
\mathrm{Ext}^{1}(\mathbb{G}_{m}, \mathbb{Z}/n\mathbb{Z}) = \mathbb{Z}/n\mathbb{Z}
\]
\note{dans cette ligne le second argument est surchargé
(« $\mathbb{Z}/n\mathbb{Z}$ » sur un autre terme, peut-être ${}_{n}\mu$).}
\[
\mathrm{Ext}^{1}(\mathbb{G}_{m}, \mathbb{Z}/n\mathbb{Z}) \simeq {}_{n}\check{\mu}
\qquad
\mathrm{Ext}^{1}(\mathbb{G}_{m}, \mathbb{Q}/\mathbb{Z}) \simeq \varinjlim_{n}\ {}_{n}\check{\mu}
\]
\[
\mathrm{Hom}\bigl(T_{\cdot}(\mathbb{G}_{m}), \mathbb{Q}/\mathbb{Z}\bigr)
\]
\[
\boxed{\mathrm{Ext}^{1}_{k\text{-}\mathrm{gr}}(\mathbb{G}_{m}, \mathbb{Q}/\mathbb{Z}) \quad \ldots\ k^{*}}
\qquad \simeq \mathrm{Hom}\bigl(\mathbb{G}_{m}(k), \mathbb{Q}/\mathbb{Z}\bigr)
\]
\note{dans le cadre, entre $\mathrm{Ext}^{1}$ et $k^{*}$, deux mots
soulignés, illisibles (\ill{} \ill{}).}
\note{le cadre, en bas à droite, est ouvert à droite ; le
« $\simeq \mathrm{Hom}(\mathbb{G}_{m}(k), \mathbb{Q}/\mathbb{Z})$ » est
écrit dessous.}

\page{109}
\note{feuillet dactylographié, annoté à l'encre et au crayon, et barré de
trois longs traits verticaux au crayon. Le texte tapé est donné tel quel ;
ses corrections à la main sont rendues comme biffures et ajouts
(au crayon ou à l'encre bleue), les mots barrés à la machine (« xxx ») ne
sont pas relevés.}

IV 7 Oubli au nº 7.6.

\emph{Proposition} 6.7. Soient $f \colon X \to Y$ un morphisme
\struck{de présentation finie} \struck{et} \add{propre, avec $Y$ loc.\
noeth.}, $F$ un Module \struck{de présentation finie} \add{cohérent} sur
$X$ \add{de support $X$}. Pour tout $y \in Y$, soit $\chi(y)$ la somme des
multiplicités géométriques totales \add{sur $k(y)$} de $F_{y}$ en les points
maximaux de $\operatorname{supp} F_{y}$. Alors la fonction
\struck{$y$}~$\mapsto \chi($\struck{$y$}$)$ sur $Y$ est semi-continue
supérieurement \add{en $y$}.
\marginal{On suppose $f$ univ.\ ouvert en les pts maximaux de $X_{y}$, et
$F_{y}$ sans cycles premiers ass.\ immergés}
\note{la note marginale, à l'encre, est reliée par un trait au point final
de l'énoncé ; « propre » tapé est gardé, le « et » qui le précédait biffé ;
le $\chi$ est une lettre tapée surchargée à la main.}

Le procédé bréveté nous ramène au cas où $Y$ est noethérien. On sait (réf)
que la fonction $\chi$ est constructible, il reste donc (réf) à prouver que
$y \in \overline{y'}$ implique $\chi(y) \geqslant \chi(y')$. Le procédé
habituel nous permet alors de nous ramener au cas où $Y$ est le spectre d'un
anneau de valuation discrète, dont $y$ ($y'$) est le point fermé
(générique), et à prouver dans ce cas l'inégalité correspondante sur les
longueurs :

\emph{Lemme} 6.8. Soit $f \colon X \to Y$ un morphisme propre, avec $Y$
spectre d'un anneau de valuation discrète, $F$ un Module cohérent sur $X$
\add{de support $X$}, $F_{0}$ et $F_{1}$ les Modules induits sur la fibre
spéciale resp.\ la fibre générique, $x_{i}$ ($x'_{j}$) les points maximaux de
$\operatorname{supp} F_{0}$ ($\operatorname{supp} F_{1}$). Alors on a
\[
\textstyle\sum_{i} \operatorname{long} F_{0\,x_{i}} \geqslant \sum_{j} \operatorname{long} F_{1\,x'_{j}}
\]
Si l'inégalité est vérifiée, alors $F$ est plat sur $Y$ aux points
maximaux de $\operatorname{supp} F_{0}$\struck{, donc toute composante de
supp $F$ domine $Y$}.
\note{au-dessus de « les points maximaux de supp $F_{0}$ », une insertion à
l'encre, barrée avec la ligne : « Supposons que tt comp.\ irr.\ de $X$
domine $Y$, et $F_{0}$ \ill{} $(S_{1})$ ». Le signe tapé de l'inégalité et
le premier $\sum$ sont en partie repassés à la main ; le texte tapé porte
« $\sum$ » devant le second membre seulement.}

Soit $t$ une uniformisante pour $V = \underline{O}_{Y,y'}$, soit $N$ le
sous-Module de $F$ réunion des annulateurs des $t^{n}$, donc on a une suite
exacte
\[
0 \to N \to F \to G \to 0
\]
d'où une suite exacte (grâce au fait que $t$ est $G$-régulier) :
\[
0 \to N/tN \to F/tF \to G/tG \to 0
\]
Cela montre que les deux sommes relatives à $F$ qu'il s'agit de comparer,
s'obtiennent en ajoutant les sommes analogues pour $N$, $G$. Comme
$N_{1} = 0$, pour prouver notre inégalité, il suffit donc de prouver
l'inégalité analogue pour $G$ ; et on ne peut avoir égalité que si
$N/tN = N_{0}$ est nul aux points maximaux de $\operatorname{supp}$ \ldots

N.B. J'ignore si les conditions \ill{} \ill{} \ill{} \ill{} \ill{}
\uncertain{suffisent} pour dire $F$ plat$/Y$ aux pts maximaux de $X_{y}$.
Question liée : $A$ anneau local, noeth., $M$ module de type fini sur $A$,
$t \in \mathfrak{m}_{A}$ \struck{\ill{}} \ill{} \ill{} \uncertain{non}
\ill{}, \ill{} \ill{} \ill{} $M = {}$\ill{}$\,A$, \ill{} $A/t_{0} = {}$\ill{},
\ill{} \ill{} $m \in \operatorname{Ass} M/tM$ !
\note{le N.B., à l'encre, occupe le bas de la page, en écriture très
rapide ; seules les formules et quelques mots sont lus.}

\page{111}
\note{au crayon ; feuille partagée par un trait vertical et deux traits
horizontaux en six cases. Dans chaque case, un groupe, une flèche ondulée
vers une flèche de gerbes, et, dans une ellipse, deux valeurs repérées
par les degrés $0$ et $1$ (au-dessus, entre crochets). $\mathrm{Grbg}$ rend
son abréviation, lue comme « gerbe(s) ».}

\begin{itemize}
\item $\mu_{n}$ \quad \struck{\ill{}} $(n, p) = 1$ \quad $\leadsto$ \quad
$[\,\mathrm{Grbg}(\mathbb{G}_{m}) \xrightarrow{\ n\ } \mathrm{Grbg}(\mathbb{G}_{m})\,]$ ;
$[\,0 \quad 1\,]$ : $\mu_{n}(k)$, $\mathbb{Z}/n\mathbb{Z}$.
\item $\mathbb{Z}/n\mathbb{Z}$ \quad $(n, p) = 1$ \quad $\leadsto$ \quad
\struck{\ill{}} ; $0$, $1$ : $\mathbb{Z}/n\mathbb{Z}$, $\check{\mu}_{n}$.
\item $\mu_{p}$ \quad $\leadsto$ \quad
$\mathrm{Grbg}(\mathbb{G}_{m}) \xrightarrow{\ p\ } \mathrm{Grbg}(\mathbb{G}_{m})$ ;
$[\,0 \quad 1\,]$ : $0$, $K^{*}/K^{*p}$.
\item $\mathbb{Z}/p\mathbb{Z}$ \quad $\leadsto$ \quad
$\mathrm{Grbg}(\mathbb{G}_{a}) \xrightarrow{\ \wp\ } \mathrm{Grbg}(\mathbb{G}_{a})$ ;
$\mathbb{Z}/p\mathbb{Z}$, $H^{1}(K, \mathbb{Z}/p\mathbb{Z}) \simeq K^{+}/\wp K^{+}$.
\item $\alpha_{p}$ \quad $\leadsto$ \quad
$\mathrm{Grbg}(\mathbb{G}_{a}) \xrightarrow{\ \mathrm{Frob}\ } \mathrm{Grbg}(\mathbb{G}_{a})$ ;
$0$, $K^{+}/\wp K^{+p}$ ($\wp$ incertain).
\item $\alpha_{p}$ \quad $\leadsto$ \quad \ill{} ;
$0$, $K^{+}/\wp K^{+p}$ ($\wp$ incertain).
\item $\mathbb{G}_{m}$ \quad ?? \quad : $K^{*}$, $0$.
\item $\mathbb{Z}$ : $\mathbb{Z}$, $0$ \quad ?? \quad $H^{1}(K, \mathbb{Q}/\mathbb{Z})$
\end{itemize}
\note{case $\mathbb{Z}/n\mathbb{Z}$ : la flèche de gerbes est raturée
(on devine $\mathbb{Q}/\mathbb{Z} \to \mathbb{Q}/\mathbb{Z}$). Case
$\mathbb{Z}/p\mathbb{Z}$ : sous $H^{1}(K, \mathbb{Z}/p\mathbb{Z})$, la
glose « cf.\ corps de classes local » ; plus bas, hors de l'ellipse,
« $\mathrm{Ext}^{0}(K^{*}/K^{*p}, \mathbb{Q}/\mathbb{Z})$ » et
« $\mathrm{Ext}^{1}(K^{*}/K^{*p}, \mathbb{Q}/\mathbb{Z})$ », le second
marqué « $\wr$ ». Le premier quotient de la case $\mathbb{Z}/p\mathbb{Z}$
est d'abord écrit $K^{*}/pK^{*}$ ; dans les cases $\alpha_{p}$, le
symbole devant $K^{+p}$ est surchargé. En bas à gauche, un $\mathbb{Z}$ et
« ?? » hors des ellipses ; la case de droite du bas n'a pas de flèche de
gerbes. Un trait ondulé ferme la page.}

\page{112}
\note{à l'encre bleue, écriture très rapide ; la page commence au milieu
d'une phrase et est barrée de deux longs traits obliques au crayon.
Beaucoup de passages sont soulignés ou barrés ; seules les formules et une
partie des mots sont lus.}

\ill{}, (\struck{\ill{}} \ill{} \struck{\ill{}} $\lambda_{x} \in
\gamma(\xi, \mathfrak{X})$ (\struck{\ill{}} devant $\xi$) \add{par tt $x \in X$,}) $=
\operatorname{Is}_{K}\bigl(K(\mathcal{O}_{\xi}), K(\mathcal{O}_{x})\bigr) =
\operatorname{Is}_{K}\bigl(K', K(\mathcal{O}_{x})\bigr)$, i.e.\ \ill{} \ill{}
\ill{} \ill{} \uncertain{unifié} \struck{\ill{}} \ill{} $K'$ \ill{} \ill{}
\ill{} $x$. \ill{} \ill{} une famille d'homomorphismes
\[
\varphi_{x} \colon \mathrm{Gal}(\overline{K}_{x}/K_{x}) \xrightarrow{\ \text{épim.}\ } \mathrm{Gal}(K'/K^{D}_{x}) \subset \mathrm{Gal}(K'/K)
\]
(homomorphismes de réciprocité)\marginal{applications de réciprocité i.e.\
\ill{} \ill{} \ill{}}. \struck{\ill{}} \struck{\ill{}} \ill{} \ill{}
(\add{invariant par $\mathrm{Gal}(K'/K)$}) \ill{} \ill{} $\operatorname{Im}
\varphi_{x}$ \ill{} \ill{} \struck{\ill{} \ill{}} \add{\ill{}} \ill{}
$\mathrm{Gal}(K'/K^{D,\mathrm{par}}_{X})$, \ill{} $K^{D,\mathrm{par}}_{X}$
\ill{} le plus grand \ill{} \ill{} extension (\add{galoisienne}) \ill{} $K'$,
\ill{} \ill{} \ill{} \uncertain{complètement} \ill{} \ill{} \ill{}
$x \in X$ \struck{\ill{}} \uncertain{décomposée} \ill{} \ill{}.
\note{le $\xi$ et le $\mathfrak{X}$ de la première ligne sont indicés d'un
$\xi$ minuscule ; les exposants $D$ de $K^{D}_{x}$ et de $K^{D,\mathrm{par}}$
sont surchargés.}

\ill{} \ill{}, \ill{} \ill{} \ill{} \ill{} \ill{} \ill{} les \ill{} \ill{}
\[
\mathrm{Gal}(\overline{K}_{x}/\overline{K}^{\,i}_{x}) \xrightarrow{\ \text{épim.}\ } \mathrm{Gal}(K'/K^{i}_{x}) \subset \mathrm{Gal}(K'/K)
\]
\ill{} \ill{} \ill{} \ill{} (\add{invariant par \ill{}}) \ill{} \ill{}
$\varphi_{x}\bigl(\mathrm{Gal}(\overline{K}_{x}/K^{I}_{x})\bigr)$ le groupe de
Galois $\mathrm{Gal}(K'/K^{I}_{X})$, \ill{} $K^{I}_{X}$ \ill{} la plus
grande \ill{} extension \struck{\ill{}} \ill{} de $K'$ non ramifiée \ill{}
$X$.
\note{les exposants $I$ de $K^{I}_{X}$ sont écrits sur un mot biffé
(« \ill{} »), de même, sous le second, un « \struck{\ill{}} ».}

\page{113}
\note{au crayon, puis à l'encre pour la dernière ligne. En haut, isolés :
« $\check{\mu}_{n}$ », « $\mathbb{Z}/p\mathbb{Z}$ », « $\mathrm{Ext}$ ».}
\[
\boxed{H^{i}\bigl(K, D(G)\bigr) \simeq \mathrm{Ext}^{1-i}\bigl(f_{*}(G), \mathbb{Q}/\mathbb{Z}\bigr)}
\]
\note{le premier $K$ est écrit sur une autre lettre ; sous la ligne,
isolé, « $A(k)$ ».}
\[
\underline{\mathrm{Ext}}^{i}( \qquad
\mathrm{Ext}^{0}\bigl(A(k), \mathbb{Q}/\mathbb{Z}\bigr)^{1-i}
\]
\[
\boxed{f_{*}D(G) \simeq \bigl(T f_{*}G\bigr)(1)}
\qquad
= \mathrm{Ext}^{i}(A, \mathbb{Q}/\mathbb{Z})(k)^{1-i}
\]
\[
H^{i}\bigl(K, D(G)\bigr) \simeq \mathrm{Ext}^{i}\bigl(f_{*}(G), \mathbb{Q}/\mathbb{Z}\bigr)
\qquad
\mathrm{Ext}^{1+h-i}(A, \mathbb{Q}/\mathbb{Z})^{1+h-i}
\]
\note{les deux formules de droite forment une colonne ; les exposants $h$
sont une lecture incertaine. Un trait relie la dernière au cadre suivant.}
\[
\boxed{H^{i}\bigl((T f_{*}G)(1)\bigr) = H^{i-1}(T f_{*}G) = \underline{\underline{\mathrm{Ext}}}^{i-1}(f_{*}G, \mathbb{Q}/\mathbb{Z})}
\]
\struck{$\boxed{H^{i}(K, D(G)) \ \ldots\ \xrightarrow{\ \sim\ } \underline{\underline{\mathrm{Ext}}}^{i-1}(f_{*}G, \mathbb{Q}/\mathbb{Z})}$}
\note{dans ce cadre biffé, entre $D(G))$ et la flèche, un terme lui-même
raturé (\ill{}).}
\[
\mathrm{Ext}^{i-1}\bigl(A(-1), \mathbb{Q}/\mathbb{Z}\bigr) = \mathrm{Ext}^{i}\bigl(M(\mathbb{Q}/\mathbb{Z})(1)\bigr)^{i-1} = \mathrm{Ext}^{i-1}\bigl(H^{1-i}(f_{*}G)\bigr)\ \ \mathrm{Ext}\ (\ldots
\]
\note{la dernière partie, « $\mathrm{Ext}^{i-1}(H^{1-i}(f_{*}G))$ » et
dessous « $\mathrm{Ext}\ ($ », est écrite au-dessus de la ligne, à droite ;
le $M$ est une lecture incertaine.}
\[
\boxed{H^{*}(K, DG) \Longleftarrow \mathrm{Ext}^{p}\bigl(H^{-q}(K, G), \mathbb{Q}/\mathbb{Z}\bigr)}
\]
\[
0 \to \mathrm{Ext}^{1}_{k\text{-}\mathrm{gr}}\bigl(H^{1-i}(K, G), \mathbb{Q}/\mathbb{Z}\bigr) \to H^{i-1}(K, DG) \to \mathrm{Hom}_{k\text{-}\mathrm{gr}}\bigl(H^{-i}(K, G), \mathbb{Q}/\mathbb{Z}\bigr)
\]
\note{à l'encre ; les exposants $1-i$, $i-1$ et $-i$ sont lus tels quels,
sans les accorder. Le cadre précédent porte à gauche un trait vertical.}

\page{114}
\note{feuillet dactylographié paginé « III-49 », scanné tête-bêche, verso
de la page 113 et sans rapport direct avec elle ; texte d'une rédaction du
chapitre III des \emph{Éléments}, non recomposé. Il porte la fin d'une
démonstration de finitude pour un morphisme propre (dévissage, lemme de
Chow, $X' \to X$ projectif birationnel, $F = g_{*}(\mathcal{O}_{X'}(n))$,
suite spectrale de Leray) et le début du Corollaire 1 (« Soit $Y$ un
préschéma localement noethérien. Pour tout morphisme propre $f \colon X \to
Y$, l'image directe par $f$ de \ldots »). Interventions à la main, de main
non établie :
\begin{itemize}
  \item « projectif » inséré à l'encre bleue avant « surjectif », et « et »
    avant « birationnel » : « un morphisme projectif surjectif et
    birationnel $g$ » ;
  \item le passage « et $f$ séparé, on en déduit (II,4,6,th.3 (v)) que $g$
    est aussi un morphisme projectif ; comme $X$ est noethérien » est
    raturé à l'encre bleue ;
  \item la référence « (FAC,I,2,12,prop.1) » est ajoutée à l'encre bleue
    après « le support de $F$ contient un ouvert partout dense de $X$ » ;
  \item « même raisonnement » est souligné d'un trait ondulé, avec une
    croix en marge ; « un » (« un $\mathcal{O}_{X}$-Module ») est souligné,
    « $K'$ » souligné de rouge ; plusieurs lettres ($\mathcal{O}_{X}$, $G$,
    $S$, $S_{1}$, $F$) sont surlignées de jaune.
\end{itemize}}

\page{115}
\note{à l'encre bleue ; feuillet déchiré à gauche.}
\[
H^{*}\bigl(\Delta^{2}T(A)\bigr) \Longleftarrow (H^{p}\Delta^{2})\bigl(H^{q}(K, A)\bigr)
\]
\[
E^{pq}_{2} = \ \Bigl\lbrace\ 0 \text{ si } q \neq 0,\ p = 1
\]
\note{entre « $=$ » et l'accolade, un terme raturé (\struck{\ill{}}).}
\note{à gauche, un petit diagramme des $E_{2}^{pq}$ : axes $p$, $q$, trois
points marqués aux positions $(-1,0)$, $(0,0)$, $(0,1)$, joints par un
carré et une diagonale.}
\[
\begin{cases}
(H^{-1}\Delta^{2})\bigl(H^{0}(K, A)\bigr) \subset \pi_{1}(A) \\
(H^{0}\Delta^{2})\bigl(H^{0}(K, A)\bigr) \\
(H^{0}\Delta^{2})\bigl(H^{1}(K, A)\bigr) = H^{1}(K, A)
\end{cases}
\]
\note{à droite de la première ligne, la glose soulignée
« \uncertain{composante} » ; à droite de la deuxième, sur deux lignes et
entourée, « \struck{\ill{}} est de torsion \ill{} \ill{} \ill{} », qui
remplace un mot biffé.}

$\Delta T(P_{\cdot})$ et $T(P'^{\cdot})$ ont \uncertain{même}
$\Delta$.
\note{la fin de la phrase est soulignée deux fois ; $P$ est un $P$
cursif.}
\[
\boxed{\underline{\underline{\mathrm{Ext}}}^{\circ}(P, \mathbb{Q}/\mathbb{Z}) = 0}
\]
$\mathrm{Ext}^{p}(H$
\[
\Bigl|\ T(P') \Longleftrightarrow \Delta T(P_{\cdot})
\]

\struck{\emph{Remarque} 2.7. J'ignore si les conditions
d'équidimensionnalité de 1.1 (iii) sont essentielles ; de même dans 2.6.

\emph{Prop.} 2.8. Constructibilité des Cohen-Macaulay.}
\note{ces lignes, au bas de la page, sont écrites tête-bêche par rapport au
reste et barrées de trois traits obliques ; elles n'appartiennent pas au
calcul qui précède.}

\page{116}
\note{à l'encre bleue ; demi-feuillet.}

$A'_{0} = A'$ \qquad $\pi_{0}(A)$ et $\pi_{0}(B)$ duaux
\[
(3)\ \begin{cases}
H^{1}(K, A) \simeq \underline{\mathrm{Ext}}^{1}_{k\text{-}\mathrm{gr}}(B, \mathbb{Q}/\mathbb{Z}) \simeq \underline{\mathrm{Ext}}^{1}_{k\text{-}\mathrm{gr}}(B'^{\circ}, \mathbb{Q}/\mathbb{Z}) \times \underline{\mathrm{Ext}}^{1}_{k\text{-}\mathrm{gr}}(V, \mathbb{Q}/\mathbb{Z}) \\
H^{1}(K, B) \simeq \underline{\mathrm{Ext}}^{1}_{k\text{-}\mathrm{gr}}(A, \mathbb{Q}/\mathbb{Z}) \simeq \underline{\mathrm{Ext}}^{1}_{k\text{-}\mathrm{gr}}(A'^{\circ}, \mathbb{Q}/\mathbb{Z}) \times \underline{\mathrm{Ext}}^{1}_{k\text{-}\mathrm{gr}}(U, \mathbb{Q}/\mathbb{Z})
\end{cases}
\]

(2) \quad $\pi_{0}(A)$ et $\pi_{0}(B)$ duaux (à valeurs dans
$\mathbb{Q}/\mathbb{Z}$)

(1) \quad $0 \to \pi_{1}(A'_{0})$
\note{les trois lignes (3), (2), (1) sont encadrées ; en marge droite du
cadre, un mot écrit en long, raturé (\struck{\ill{}}).}
\[
0 \to U \to A \to A' \to 0
\]
\note{sous $U$ : « unipot.\ connexe » ; sous $A'$ : « \struck{\ill{}}
\uncertain{semi-ab.} connexe ».}
\[
\mu(A) = 2\alpha(A^{\circ}) + \mu(A^{\circ})
\qquad
0 \to \mathrm{Ext}^{1}(A'^{\circ}, \mathbb{Q}/\mathbb{Z}) \to \mathrm{Ext}^{1}(A, \mathbb{Q}/\mathbb{Z}) \to \underline{\mathrm{Ext}}^{1}(U, \mathbb{Q}/\mathbb{Z}) \to 0
\]
\[
(\mathbb{Q}/\mathbb{Z})^{\mu(A)} \qquad \mathrm{Ext}^{1}(A, \mathbb{Q}/\mathbb{Z}) \simeq H^{1}(K, B)
\]
\note{sous $\alpha(A^{\circ})$ et $\mu(A^{\circ})$, soulignés, deux gloses
serrées, l'une raturée : « \ill{} \struck{\ill{}} \ill{} » et
« \uncertain{multiplicatif} » ; $(\mathbb{Q}/\mathbb{Z})^{\mu(A)}$ est écrit
sous le premier $\mathrm{Ext}^{1}$ (le $\mathbb{Q}/\mathbb{Z}$ surchargé),
$H^{1}(K, B)$ sous le second, relié par « $\wr$ » ; sous
$\underline{\mathrm{Ext}}^{1}(U, \mathbb{Q}/\mathbb{Z})$, une glose sur
deux lignes : « \ill{} \uncertain{groupe} \ill{} \ill{} \ill{} ».}

\page{117}
\note{à l'encre, liste numérotée 1.--15., coupée par trois longs traits
verticaux au crayon. Une accolade ondulée réunit 4.--5. et 6.--7., un
crochet 8.--10., une accolade 11.--12. ; un trait noir épais marque 13.}

\emph{Morphismes plats etc.}
\begin{enumerate}
\item[1.] Changement de base, \struck{platitude}
\item[2.] $\varprojlim$ de préschémas, propriétés constructibles
\item[3.] (P.) Morphismes réguliers etc. Applications à l'algèbre locale.

Pourrait passer en Ch.\ 0 ?
\item[4.] Morphismes plats de présentation finie : critères de platitude
\item[5.] Immersions régulières etc
\item[6.] Dimension, morphismes équidimensionnels
\item[7.] Morphismes univt.\ ouverts.
\item[8.] \ill{} \uncertain{diff.}. Morphismes
\uncertain{diff.} simples.
\item[9.] Morphismes simples
\item[10.] Morphismes étales, morphismes non ramifiés
\item[11.] Fibres d'un morphisme plat de présentation finie. Morphismes
séparables, normaux etc
\item[12.] Sections hyperplanes, multisections
\item[13.] Prolongements infinitésimaux
\item[14.] Préschémas en groupes. \uncertain{Espaces} \uncertain{algébriques}
\uncertain{de} \uncertain{Serre}.
\item[15.] Formes différentielles et dualité
\end{enumerate}
\note{« Morphismes plats etc. », en tête, est encadré d'un arc qui descend
vers 1. Au bas de la page, tête-bêche et barrée de grandes croix, une liste
de choses à faire, en anglais et sans mathématique, qu'on ne transcrit pas ;
elle porte « September 1963 ».}

\page{118}
\note{au crayon.}
\[
\underline{\mathrm{Ext}}^{i}(A, \mu_{\infty}) = \lbrace 0,\ {}_{\infty}B,\ 0,\ 0
\]
\[
\underline{\underline{\mathrm{Ext}}}^{i}(A, \mu_{\infty}) \simeq H^{i-1}(K, {}_{\infty}B) = \lbrace H^{0}(K, {}_{\infty}B),\ H^{1}(K, {}_{\infty}B)^{\times},\ 0
\]
\[
H^{i}(K, A) = \lbrace A(K)^{\times},\ H^{1}(K, A),\ 0,\ 0,\ 0
\]
\struck{$B \quad H^{1}(K, {}_{\infty}B) \to H^{1}(K, B) \to 0$}
\note{les valeurs sont en colonne à droite d'accolades ; la première
colonne de la deuxième ligne commence par un « $0$ » écrit au-dessus. Le
petit signe en exposant après $H^{1}(K, {}_{\infty}B)$ et $A(K)$ est un
$\times$ ou une marque de renvoi. Le troisième terme de la colonne de
$H^{i}(K, A)$ est surchargé.}
\[
\to H^{0}(K, B) \otimes \mathbb{Q}/\mathbb{Z} \to H^{1}(K, {}_{\infty}B) \to H^{1}(K, B) \to 0
\]
\note{le premier terme est écrit sur un autre ($H^{0}(K,B)$ surchargé).}

\struck{$H^{1}(K, A) \to$}

\begin{tikzcd}[column sep=small, row sep=small, nodes={font=\scriptsize}]
0 \arrow[r] & H^{0}(K, {}_{\infty}B) \arrow[r] & DH^{1}(K, A) \arrow[d, no head, "\parallel"'] \arrow[dr, bend left=20] & & \\
& & H^{0}(K, B) & H^{0}(K, B) \otimes \mathbb{Q}/\mathbb{Z} \arrow[d, bend left=20] & \\
& & & H^{1}(K, {}_{\infty}B) \arrow[r] & DH^{0}(K, A) \arrow[r] & 0
\end{tikzcd}

\note{redessiné. La première ligne continuait par
« \struck{$\to H^{1}(K, {}_{\infty}B)$} » ; une flèche courbe part de
$DH^{1}(K, A)$ vers $H^{0}(K, B) \otimes \mathbb{Q}/\mathbb{Z}$, une autre
de là vers $H^{1}(K, {}_{\infty}B)$. Sous $DH^{0}(K, A)$ : « $\parallel$
$H^{1}(K, B)$ ». Le tout est entouré d'un trait qui l'isole.}
\[
0 \to H^{0}(K, {}_{\infty}B) \to DH^{1}(K, A) \to Q \to 0
\]
\[
0 \to N \to H^{1}(K, {}_{\infty}B) \to DH^{0}(K, A) \to 0
\]
$Q = $ \struck{$H^{1}(K, B)$} \struck{$\operatorname{Ker} H^{0}(K, B) \to H^{0}(K, B) \otimes \mathbb{Q}/\mathbb{Z})$}
\[
Q \subset N \qquad H^{0}(K, B) \to H^{0}(K, B) \otimes \mathbb{Q}/\mathbb{Z}
\]

\page{119}
\note{au crayon. En haut à droite, isolés : « $\mathbb{Z}/n\mathbb{Z}$ »,
« $B$ ».}
\[
\underline{\mathrm{Ext}}^{i}(A, {}_{n}\mu) = \lbrace 0,\ {}_{n}B,\ 0,\ 0
\]
\[
\underline{\underline{\mathrm{Ext}}}^{i}(A, {}_{n}\mu) \simeq H^{i-1}(K, {}_{n}B)
\simeq \mathrm{Ext}^{i-1}\bigl(T(A), \mathbb{Z}/n\mathbb{Z}\bigr)
\]
\note{le dernier $\simeq$ est écrit verticalement.}
\[
H^{i}(K, {}_{n}B) \xrightarrow{\ \sim\ } \mathrm{Ext}^{i}_{k\text{-}\mathrm{gr}}\bigl(T(A), \mathbb{Z}/n\mathbb{Z}\bigr)
\]
\[
H^{i}(K, {}_{\infty}B) \simeq \mathrm{Ext}^{i}_{k\text{-}\mathrm{gr}}\bigl(T(A), \mathbb{Q}/\mathbb{Z}\bigr)
\]
\note{sous $T(A)$, une flèche double verticale descend vers « $A$ »,
et à côté « $H^{1}(K, A)$ » (le $A$ repassé), sous une accolade
horizontale : « $V \times$ \struck{\ill{}} $D$ -- divisib. ».}

${}_{n}B(K)$, $H^{1}(K, {}_{n}B)$, $0$ \ldots \qquad
[dual de $H^{1}(K, {}_{n}A)$]
\note{à gauche, un diagramme de suite spectrale : axes $p$, $q$, deux
points sur la ligne $q = 0$ ; à la colonne $p = 0$,
« $\Delta_{n}(\pi_{0}(A))$ » et « \struck{\ill{}}
$\Delta_{n}(H^{1}(K,A)')$ » ; dans une ellipse,
« $\mathrm{Ext}^{1}(A, \mathbb{Z}/n\mathbb{Z})$ ».}
\[
\Delta\bigl(H^{1}(K, A)'\bigr)_{n} = \angle\bigl({}_{n}H^{1}(K, A)\bigr)
\]

\begin{enumerate}
\item[1°)] la partie divisible de $H^{1}(K, A)$ est \emph{divisible}
\item[2°)]
\end{enumerate}
\note{« divisible », en fin de 1°), est souligné trois fois ; la lecture
de toute la phrase est incertaine.}
\[
0 \to \mathbb{Z}/n\mathbb{Z} \to \mathbb{Q}/\mathbb{Z} \xrightarrow{\ n\ } \mathbb{Q}/\mathbb{Z} \to 0
\]
\begin{enumerate}
\item[(i)] $H^{1}(K, A)'$ divisible, i.e.\ $\angle H^{1}(K, A)'$ sans
torsion \struck{$\mathbb{Z}$}
\item[(ii)] $H^{1}(K, {}_{n}B) \xrightarrow{\ \sim\ } \mathrm{Ext}^{1}_{k\text{-}\mathrm{gr}}(A, \mathbb{Z}/n\mathbb{Z})$ \quad ?
\item[(iii)] $n$
\end{enumerate}
\note{« sans torsion » est écrit au-dessus d'un $\mathbb{Z}$ biffé ; le
signe $\angle$ est rendu tel qu'il est tracé.}

\page{120}
\note{feuillet dactylographié paginé « III-42 », verso de la page 119 et
sans rapport direct avec elle ; texte d'une rédaction du chapitre III des
\emph{Éléments}, non recomposé : la caractéristique d'Euler-Poincaré
$\chi_{A}(F) = \sum (-1)^{i} \operatorname{long}(H^{i}(X, F))$ (15), son
additivité (Proposition 13, (16)), le genre arithmétique, et le
Théorème 2 (i)--(iv) sur la fonction $n \mapsto \chi_{A}(F(n))$ (polynôme en
$n$, égal à $\operatorname{long}\Gamma(X, F(n))$ pour $n$ assez grand,
coefficient dominant $\geqslant 0$, degré égal à la dimension du support),
avec le début de la démonstration. Interventions à la main, de main non
établie :
\begin{itemize}
  \item « nombre entier » : « entier » séparé de « nombre » par un trait ;
  \item un « o » au-dessus de « on l'appelle la », relié à une note en
    marge gauche, « \uncertain{si} $A$ est local » ;
  \item « Soit $F$ un $\mathcal{O}_{X}$-Module cohérent. » inséré avant
    « Théorème 2 » par un crochet ; dans (i), « est égale à un polynôme en
    $n$ » (« égale à » et « en $n$ » ajoutés) ;
  \item en marge, un « $>$ » en face de (iii), et, le long du dernier
    paragraphe, un trait vertical avec « \uncertain{lourde} » ;
  \item une lettre ajoutée au-dessus de « exacte » ; plusieurs lettres
    surlignées de jaune.
\end{itemize}}

\end{document}
